\documentclass[11pt]{article}

\usepackage[margin=1.05in]{geometry}
\usepackage{amsmath,amssymb,amsthm}
\usepackage{booktabs}
\usepackage{enumitem}
\usepackage{xcolor}
\usepackage[ruled,vlined,linesnumbered]{algorithm2e}
\usepackage[colorlinks=true,linkcolor=blue!45!black,citecolor=blue!45!black,urlcolor=blue!45!black]{hyperref}

\newcommand{\R}{\mathbb{R}}
\newcommand{\E}{\mathbb{E}}
\newcommand{\tr}{\operatorname{tr}}
\newcommand{\xmap}{x_{\star}}
\newcommand{\code}[1]{\texttt{\small #1}}
\newcommand{\Rhat}{\widehat{R}}
\newcommand{\khat}{\widehat{k}}

\theoremstyle{plain}
\newtheorem{proposition}{Proposition}
\theoremstyle{definition}
\newtheorem{condition}{Condition}
\theoremstyle{remark}
\newtheorem{remark}{Remark}

\title{\bfseries Profiling the states out:\\[2pt]
posterior inference for MAGI without a Gaussian assumption}
\author{}
\date{\today}

\begin{document}
\maketitle

\begin{abstract}
\noindent
MAGI \cite{yang2021} infers ODE parameters by placing a Gaussian process prior on the state and
constraining it to satisfy the dynamics on a discretisation grid. Its posterior is dominated
numerically by the states, which outnumber the parameters by two orders of magnitude, and the
standard summary---a Gaussian at the maximum a posteriori point---is mis-centred by a full
posterior standard deviation on a benchmark as benign as FitzHugh--Nagumo.

We argue that the states should not be approximated at all. They are the part of the problem the
data and the GP prior constrain most tightly, and integrating them out by Laplace at each parameter
value leaves a $p$-dimensional integral with $p$ between $3$ and $7$, which can be done directly.
The resulting posterior is a mixture, one Gaussian in the states per parameter node, and makes no
Gaussian assumption about the parameters whatsoever. On every benchmark in our suite that has a usable
reference, its largest parameter error is at or below the level at which two independent halves of
that reference agree with each other: $0.0126$ against a floor of $0.0100$ on FitzHugh--Nagumo,
$0.0093$ against $0.0081$ on HIV, and $0.0119$ against $0.0405$ on the chaotic Lorenz system, where
the mode alone is $1.80$ out. It costs between three and twenty-three seconds against roughly two
hundred and up for a reference chain, and it reports its own reliability without one.

Getting there required correcting several things upstream, of which the largest was not in the
inference at all. MAGI's GP hyperparameters are fitted by an optimiser started at
$\phi_1 = \phi_2 = \sigma = 1$ regardless of the data's units; on a benchmark whose states have
marginal variance $5.6\times10^4$ this settles on a lengthscale four orders below the grid
spacing, so the GP imposes no smoothness, the trajectory interpolates the observations and is
unconstrained between them, and the parameters absorb the difference. Putting that fit on a
scale-free footing moves the trajectory error on that system from $69.5\%$ to $0.4\%$ and takes its
Hessian condition number from $4\times10^{17}$ to $4\times10^{2}$. Separately, the Gauss--Newton
normal equations square the Jacobian's condition number, which for parameters in mixed units is
repaired by one line of diagonal scaling; the benchmark suite's forward Euler integrator carried
more error than the observation noise it was competing with; and MAGI's implementation omits the
prior $\pi_\Theta(\theta)$ that its own derivation includes.

We therefore also report a pre-inference diagnosis, costing seconds and no sampling, that
establishes whether the mode is a mode, whether it is the only one, which parameters the data
determines, whether the posterior is proper, and whether the GP constrains the states at all. The
last of these is the one that mattered: the others reported the broken hyperparameter fit
faithfully as a property of the model.
\end{abstract}

\section{Introduction}

Estimating the parameters $\theta \in \R^p$ of $\dot x = f(x,\theta)$ from sparse noisy
observations is awkward because the likelihood contains an ODE solve. MAGI \cite{yang2021} removes
it: a Gaussian process prior on each state component, a discretisation grid $I$, and a constraint
that the GP satisfy the dynamics on that grid. The unknown becomes
\[
  x = (\theta, X) \in \R^d, \qquad d = p + nD,
\]
and on the four systems we study $d$ runs from $106$ to $608$ while $p$ runs from $3$ to $7$.

That imbalance is the whole story of this paper. Almost all of the dimension is state, and the
state is the part of the problem that the observations and the GP prior pin down; the parameters
are few and are what the exercise is for. Approximating everything at once with a Gaussian at the
mode---the default summary---spends its accuracy in the wrong place, and we show it is mis-centred
by a full posterior standard deviation on the parameters even when the covariance is nearly right,
so nothing in the output looks wrong. Sampling the joint posterior properly is possible but costs
minutes and, as we document, can fail silently on exactly the problems where it is most needed.

\paragraph{Contributions.}
\begin{enumerate}[leftmargin=1.5em,itemsep=1pt]
\item A posterior approximation that profiles the states out and integrates the parameters
directly (Section~\ref{sec:profiled}), reaching the reference chain's own agreement floor on every
benchmark that has a usable reference, in three to twenty-three seconds, and reporting its
reliability without a reference.
\item A Gaussian-process hyperparameter fit that is scale-free (Section~\ref{sec:gp}). MAGI's
optimiser starts at $\phi_1=\phi_2=\sigma=1$ whatever the data's units, and on states with
marginal variance $5.6\times10^4$ it converges to a lengthscale four orders below the grid
spacing, so the prior imposes no smoothness at all. This was the largest error we found, it was
upstream of everything else, and it invalidated several of our own earlier conclusions.
\item A pre-inference diagnosis built from the mode and its exact Hessian
(Section~\ref{sec:diagnosis}): mode validity, globality, identifiability, properness, whether the
GP constrains the states, and how badly conditioned the geometry a sampler would face is---in
seconds and without sampling.
\item The observation that MAGI's implementation omits the $\pi_\Theta(\theta)$ term of its own
derivation (Section~\ref{sec:prior}), together with a retraction of the demonstration we had built
on it, which turned out to be measuring the hyperparameter bug of Section~\ref{sec:gp}.
\item Two corrections without which none of the above is measurable: diagonal scaling of the
Gauss--Newton normal equations (Section~\ref{sec:solver}), and a benchmark integrator whose error
is below the observation noise (Section~\ref{sec:bench}).
\item Negative results, including that the local curvature of the profiled marginal is provably
equal to the joint Laplace's, so profiling buys nothing near the mode
(Section~\ref{sec:negative}).
\end{enumerate}

This paper continues \cite{companion}, which established the exact least-squares structure of the
MAGI posterior and a Gauss--Newton mode solver, and proposed a third-order correction to the
Laplace mean. That correction is superseded here: it helped on one of the four systems and harmed
two, and the profiled posterior of Section~\ref{sec:profiled} dominates it everywhere. Those
comparisons were made before the hyperparameter fit, so they describe a different set of
posteriors; we have not re-run them, because the correction is not the recommended method under
either.

\section{The MAGI posterior and its mode}
\label{sec:solver}

\subsection{Least-squares structure}

With the observation noise held fixed, MAGI's negative log-posterior is exactly a sum of squares
\cite{companion}. Writing $L_cL_c^\top = C^{-1}$, $L_kL_k^\top = K^{-1}$, $b = \beta^{-1}$ for the
prior tempering, and $P$ for the parameter prior precision introduced in Section~\ref{sec:prior},
\begin{equation}
\label{eq:residual}
  -2\log p(x) = \|R(x)\|^2 + \text{const},\qquad
  R(x) = \begin{bmatrix}
    \sqrt{b}\, L_c^\top (X-\mu)\\
    (X-y)/\sigma \;\text{ on observed entries}\\
    \sqrt{b}\, L_k^\top\big(f(X,\theta)-\dot\mu-m(X-\mu)\big)\\
    P^{1/2}(\theta - \theta_0)
  \end{bmatrix}.
\end{equation}
Only the third block is nonlinear and only through $f$, so $\nabla^2 R_a = \sqrt b\,(L_k^\top)_a
\nabla^2 f$ pointwise, and the exact Hessian $H = J^\top J + \sum_a R_a \nabla^2 R_a$ costs one
Jacobian assembly rather than $d$ autodiff passes. The fourth block is linear, so a parameter prior
---including a full precision matrix---leaves this structure untouched.

\subsection{The normal equations square the condition number}

The solver forms $A = J^\top J$ and factorises it, which is far faster than a QR least-squares
solve and squares $\kappa(J)$. Whether that matters is not a property of the science but of the
units. Table~\ref{tab:cond} shows what happens on HIV, whose parameters
$\theta = (36, 0.108, 0.5, 1000, 3)$ are rate constants on different scales: the diagonal of $A$
spans $5\times10^{-2}$ to $7\times10^{5}$ on the $\theta$ block alone, $\kappa(A)$ reaches
$5.6\times10^{10}$, and one symmetric diagonal scaling takes it to $4.0\times10^{2}$.

\begin{table}[t]
\centering\small
\begin{tabular}{lrrrr}
\toprule
& $\kappa(A)$ & $\kappa(DAD)$ & gain & $\operatorname{diag}(A)$ on the $\theta$ block\\
\midrule
fn     & 4.4e3 & 3.9e3 & $1.2\times$ & 3e3 -- 5e3\\
lorenz & 1.1e5 & 5.0e3 & $23\times$ & 8e1 -- 1e3\\
hes1   & 2.1e9 & 1.8e5 & $1.2\times10^{4}$ & 1e2 -- 3e7\\
hiv    & \textbf{5.6e10} & \textbf{4.0e2} & $\mathbf{1.4\times10^{8}}$ & 5e-2 -- 7e5\\
\bottomrule
\end{tabular}
\caption{Symmetric diagonal scaling of the Gauss--Newton normal equations, at the mode, in double
precision. The gain tracks the spread of the diagonal, as it must: eight orders on HIV, whose
parameters are rate constants in mixed units, and essentially nothing on FitzHugh--Nagumo, whose
three are all $O(1)$. These figures are from after the hyperparameter fit of
Section~\ref{sec:gp}; before it HIV read $\kappa(A) = 4\times10^{17}$, which we had presented as
the motivating case. The motivation is weaker than we thought and the conclusion is unchanged---the
scaling costs one line and is never harmful.}
\label{tab:cond}
\end{table}

The repair is symmetric diagonal scaling,
\begin{equation}
\label{eq:jacobi}
  D = \operatorname{diag}(A)^{-1/2},\qquad (DAD + \lambda I)\,s = -Dg,\qquad \delta = Ds,
\end{equation}
identical in exact arithmetic and completely different in floating point. Three details cost more
effort than the change. First, the existing damping was a uniform ridge
$\lambda\,\tr(A)/d\cdot I$, not Marquardt's $\lambda\operatorname{diag}(A)$; on a unit-diagonal
matrix $\lambda I$ \emph{is} the latter, so scaling makes the damping per-coordinate as a side
effect. Second, $\lambda$ becomes relative, so the previous absolute floor of $10^{-12}$ sits at
the smallest eigenvalue of the scaled matrix and never stops perturbing the solve; machine epsilon
is the meaningful floor. Third, and dominant, the \emph{initial} $\lambda$ matters: at $10^{-8}$
HIV stalls at $1.5\times10^{-3}$, while at $10^{-10}$ or below it reaches $10^{-7}$, and the other
three systems are indifferent across six orders. A diagonal floor set relative to the largest entry
is tempting and wrong---HIV's legitimate curvature range spans $10^{16}$, so any relative floor
large enough to act clips real curvature.

\subsection{Convergence cannot be tested on the gradient norm}
\label{sec:decrement}

The same reasoning applies to stopping the iteration. $\|\nabla\log p\|$ carries the units of the
log-density and of $\theta$, so no absolute tolerance transfers between problems, and its
attainable floor is set by the working precision: in float32 it bottoms out near $10^{-2}$ on
FitzHugh--Nagumo, where the default tolerance of $10^{-6}$ is unreachable by construction and the
solver runs every iteration of its budget for nothing.

The Newton decrement is the scale-free replacement and costs nothing, since $\delta = A^{-1}g$ is
already formed at each step:
\begin{equation}
\label{eq:decrement}
  \lambda_{\mathrm{N}} = \sqrt{g^\top A^{-1} g} = \sqrt{g^\top \delta}.
\end{equation}
To second order this is the remaining distance to the mode in posterior standard deviations, so a
tolerance on it means the same thing on every problem and in every precision. A tolerance of
$10^{-3}$---locate the mode to a thousandth of a standard deviation---terminates every case in the
suite in $6$ to $35$ iterations where an absolute tolerance ran to the cap:

\begin{center}\small
\begin{tabular}{lrrrrrrrr}
\toprule
& \multicolumn{2}{c}{fn} & \multicolumn{2}{c}{hes1} & \multicolumn{2}{c}{hiv} & \multicolumn{2}{c}{lorenz}\\
& f64 & f32 & f64 & f32 & f64 & f32 & f64 & f32\\
\midrule
iterations & 17 & 22 & 16 & 16 & 4 & 4 & 31 & 34\\
$\lambda_{\mathrm{N}}$ & 6.3e-4 & 1.6e-3 & 2.0e-4 & 1.6e-3 & 6.7e-6 & 4.1e-4 & 8.8e-4 & 3.5e-3\\
$\|\nabla\log p\|$ & 2.6e-3 & 8.0e-3 & 1.1e-2 & 1.4e-1 & 3.7e-4 & 1.0e-2 & 5.1e-4 & 5.5e-3\\
\bottomrule
\end{tabular}
\end{center}

A scale-free tolerance still has a precision floor, so the iteration also stops on stagnation;
otherwise float32 spends its remaining budget whenever the requested tolerance happens to sit
below what the format can deliver. The table makes the case for reporting $\lambda_{\mathrm{N}}$
rather than $\|\nabla\log p\|$ in one line: read the gradient norms alone and Hes1 in float32
($1.4\times10^{-1}$) looks like the worst-converged run in the suite by more than an order of
magnitude, while its decrement, $1.6\times10^{-3}$, is unremarkable. The gradient norm is reporting the
units of Hes1's $\theta$ block, whose curvature spans $10^{2}$ to $10^{7}$; the decrement divides
them out.

Every entry here reaches the tolerance, including HIV in float32, which until recently could not
run in single precision at all: its GP covariances factored to NaN, so every residual for two of
its three components was NaN before the first step. Once the hyperparameters were fitted correctly
HIV's $C^{-1}$ had condition $1.3\times10^{9}$, comfortably inside float64's range and two orders
outside float32's.

Our first repair was to scale the Cholesky ridge with the working precision, which made the NaN go
away and was wrong. The factors \emph{define} the residual---$L_cL_c^\top = C^{-1}$ in
\eqref{eq:residual}---so a ridge large enough to rescue a float32 factorisation, $4.9\times10^{-4}$
relative, does not stabilise the solve but replaces the model. The solver then converged accurately
to the mode of something else: every float32 mode in the suite failed the gradient check by between
$0.4$ and $15$ posterior standard deviations while reporting convergence, because it had converged.

The correct repair is that the factorisation is preprocessing and its precision is not the
iteration's to choose. The two GP precision matrices are factored in float64---from snapshots taken
before anything is cast down---and only the resulting triangular factors are moved to the working
dtype. The ridge is then $9\times10^{-13}$ in both precisions, HIV factors cleanly, and every
float32 mode passes. We record the false start because the failure it produced is the kind this
paper is about: a diagnostic correctly reporting a disagreement, and the disagreement being with a
model we had altered without meaning to.

\begin{remark}[A recurring theme]
This is the fourth place in this paper where a quantity with units was being compared against an
absolute threshold, after the conditioning of the normal equations (Section~\ref{sec:solver}), the
detection of flat directions (Section~\ref{sec:diagnosis}) and the finite-difference step
(Section~\ref{sec:stencil}). Two more follow: the agreement test between restarts
(Section~\ref{sec:diagnosis}) and, worst of all, the starting point of the hyperparameter optimiser
(Section~\ref{sec:gp}). In each case the fix is to divide by the relevant scale---the diagonal of
$A$, the diagonal of $H$, the posterior standard deviation, the data's own variance---and in each
case the unscaled version worked on the benchmark it was developed on and failed elsewhere.

The Cholesky ridge just above is the instructive exception. It has exactly the same shape---a
constant calibrated against float64 being applied in float32---and scaling it to the dtype is the
obvious move and the wrong one, because that constant is not a tolerance being compared against
anything. It is a perturbation of the model. Recognising which of the two a magic number is
matters more than recognising that it is one.
\end{remark}

\section{The missing prior}
\label{sec:prior}

MAGI's posterior, as derived in \cite{yang2021}, is
\[
  p(\theta, x(I)\mid\cdot) \;\propto\; \pi_\Theta(\theta)\,
  \exp\Big\{-\tfrac12\textstyle\sum_d\big[\text{GP} + \text{observation} + \text{ODE}\big]\Big\},
\]
with ``a general prior distribution $\pi(\cdot)$ on $\theta$''. The implementation omits
$\pi_\Theta$: its \code{theta\_conf} argument reaches only the initialiser, where it pulls the
starting point toward a guess, and never enters the log-density. Every posterior therefore carries
a flat improper prior on every parameter.

By \eqref{eq:residual} a Gaussian prior is one more residual block, so restoring it preserves the
least-squares structure, contributes a constant to $J^\top J$, adds no second-derivative term, and
lifts the $\theta$ block away from singularity. We take the root from an eigendecomposition rather
than a Cholesky because $P$ is permitted to be singular, $P=0$ recovering the previous behaviour
bit for bit.

\subsection{What we thought it fixed}

We originally reported a striking demonstration here. Under the flat prior HIV's posterior was
improper in $\lambda$---the profiled log-density fell by $0.07$ nats while $\lambda$ moved from
$-11.9$ to $99{,}988$---the Hessian carried an exact null direction, importance sampling collapsed
to an effective sample size of one, and the parameters that \emph{were} well determined, $\delta$
and $N$, sat at $0.259$ and $1788$ against a truth of $0.5$ and $1000$ with posterior standard
deviations under $2\%$ of their values. Supplying a proper prior removed the null direction, made
the posterior proper, restored the effective sample size fourteenfold and moved $\delta$ and $N$ to
$0.481$ and $962$. A control with the prior centred at half the true values still recovered $0.470$
and $940$, so the prior was not supplying the answer.

All of that was measuring the hyperparameter fit of Section~\ref{sec:gp}. With $\phi$ fitted on a
scale-free footing, HIV's posterior is proper under a flat prior, has no null direction, and places
$\lambda$ at $36.3$ against a truth of $36$. The demonstration is void.

What survives is narrower but still worth stating. The implementation omits a term its own
derivation contains, and that is a discrepancy independent of any benchmark: wherever a parameter
genuinely is unidentified, a flat prior leaves the posterior improper, and the machinery for
supplying $\pi_\Theta$ costs one residual block and preserves the least-squares structure exactly.
The mechanism we described---that in a partially identified model the estimates of the identified
parameters are set by wherever the unidentified ones come to rest---remains sound as an argument,
but we no longer have an example of it on this suite, because the suite turns out to be identified.

\begin{remark}
The episode is worth keeping in view when reading the rest of this paper. Every diagnostic we had
built reported the situation correctly given the model it was handed; what none of them could do
was question the model. The observation that broke it open was not a diagnostic at all but a plot
of the fitted trajectories.
\end{remark}

\section{The hyperparameters the inference never sees}
\label{sec:gp}

MAGI fits the GP hyperparameters $\phi = (\phi_1, \phi_2)$ once, before any inference, by
maximising a marginal likelihood per state component. Nothing downstream revisits them, and
nothing in the posterior machinery can compensate for them. They turn out to dominate everything
else in this paper.

\subsection{The symptom}

The posterior trajectories are wrong on three of the four systems, and not in a way the parameter
diagnostics detect. On HIV's $T_U$, whose observation noise is $\sigma = 3.16$ and whose true range
is $168$ to $600$:

\begin{center}\small
\begin{tabular}{lr}
\toprule
 & rms$/\sigma$\\
\midrule
MAP against the data, at observed points & $0.007$\\
MAP against the truth, at observed points & $1.06$\\
MAP against the truth, at \textbf{unobserved} points & $\mathbf{74.5}$\\
\bottomrule
\end{tabular}
\end{center}

The trajectory passes exactly through every observation and falls to $9.5\times10^{-8}$ between
them. Observations occupy every second grid point, so the interleaved points are held by the GP
prior alone---and the fitted lengthscale is $1.05\times10^{-4}$ against a grid spacing of $0.1$.
The fitted posterior is not at fault: its trajectory bands match the reference chain to within a
few percent on every system. The mode is.

\subsection{The cause, and a scale-free fit}

The hyperparameters are optimised in $\log\phi$ by BFGS from $x_0 = 0$, that is from
$\phi_1 = \phi_2 = \sigma = 1$ whatever the units of the data. HIV's $T_U$ has marginal variance
$5.6\times10^4$, so the start is five orders out, and BFGS---whose step sizes and convergence
tolerances are absolute in its own variables---settles on a lengthscale that explains the
observations as white noise.

The repair is the same one as everywhere else in this paper: optimise
$\phi = s \odot e^{u}$ from $u = 0$, with $s$ the data's own marginal variance and the FFT
lengthscale estimate the code already computes. A unit step is then a factor of $e$ regardless of
units.

That alone is not enough, because the constraint is two-sided. A scale-free start pushes HIV's
$T_I$ lengthscale to $20.56$, the entire time span, at which $K^{-1}$ becomes \emph{indefinite}
(eigenvalues $-24$ to $+11$) and its Cholesky in the solver is NaN: the Matern derivative kernel
degenerates as the lengthscale approaches the span, just as it stops constraining anything when
the lengthscale falls below the grid. The fit is confined to $[2\,\delta t,\ \mathrm{span}/4]$,
both ends being properties of the discretisation rather than of the data.

Trajectory quality is governed entirely by the resulting ratio $\ell/\delta t$
(Table~\ref{tab:gp}), and $\mathtt{diagnose()}$ now reports it, because no other diagnostic can
see it.

\begin{table}[t]
\centering\small
\begin{tabular}{lcccc}
\toprule
& \multicolumn{2}{c}{before} & \multicolumn{2}{c}{after}\\
system & $\ell/\delta t$ & trajectory error & $\ell/\delta t$ & trajectory error\\
\midrule
fn     & 13, 12 & 3.5\%, 7.2\% & 13, 12 & 3.5\%, \textbf{2.6\%}\\
lorenz & 12, 9.1, 1.1 & 22.9\%, 25.3\%, 14.2\% & 12, 9.1, 17 & \textbf{8.3\%, 11.8\%, 5.3\%}\\
hes1   & 0.14, 0.15, 0.16 & 71.6\%, 73.2\%, 100.8\% & 8.0, 6.3, 3.9 & \textbf{5.8\%, 13.8\%, 30.0\%}\\
hiv    & 0.001, 0.011, 12 & 69.5\%, 6.4\%, 0.3\% & 50, 50, 8.7 & \textbf{0.4\%, 0.9\%, 0.1\%}\\
\bottomrule
\end{tabular}
\caption{Relative error of the fitted state trajectory against the truth, before and after the
hyperparameter fit is made scale free. Wherever the lengthscale exceeds a few grid spacings the
trajectory is recovered; wherever it falls below one, the states between observations are held by
nothing.}
\label{tab:gp}
\end{table}

\subsection{What this invalidates}

The consequences reach most of the quantitative claims we had made about these benchmarks. With
the hyperparameters fitted properly, \textbf{every system recovers its true parameters to within
$1.6$ posterior standard deviations}, and the only parameters still weakly determined anywhere are
three of Hes1's seven---HIV's $\lambda$ is $36.3$
against a truth of $36$ where it had been $-11.9$, and Hes1's seven parameters, which had collapsed
to $\sim\!10^{-6}$, are all within $1.4$ standard deviations of their true values. HIV's condition
number falls from $4\times10^{17}$ to $4\times10^{2}$ and its null direction disappears, so its
posterior is proper and a reference chain for it exists.

We record this rather than quietly rewriting, because the failure mode is the instructive part. The
bug sat upstream of everything under study; it produced posteriors that looked entirely plausible;
and the diagnostics of Section~\ref{sec:diagnosis}---mode validity, identifiability, properness---
reported it faithfully, as a property of the model. What exposed it was plotting the trajectories,
which none of those diagnostics examined.

\section{Diagnosis before inference}
\label{sec:diagnosis}

A fast mode makes a set of questions affordable that are usually assumed away. All of the following
requires no sampler and costs, on these four systems, between ten seconds and a minute and a half
--- most of it in the eight Hessian factorizations of the conditioning check below, which dominates
on HIV at $d=608$ and can be dropped by setting \texttt{n\_curv=0} when only identifiability is
wanted.

\paragraph{Is it a mode?} The exact Hessian's spectrum, and, if a genuinely negative direction
exists, a step along it followed by a re-solve. On all four systems there is none. Proximity to the
mode is reported as $\sqrt{g^\top H^{-1} g}$ rather than $\|\nabla\log p\|$, for the reason in
Section~\ref{sec:decrement}: on FitzHugh--Nagumo in float32 the raw norm is $1.7\times10^{-2}$,
which fails any tolerance carried over from double precision, while the scale-free measure is
$5.6\times10^{-3}$---the mode is located to within $0.6\%$ of a posterior standard deviation, and
nothing downstream is affected.

\paragraph{Is it the only one?} Re-solving from dispersed starts, and separately by continuation in
the prior tempering $\beta^{-1}$: at small $\beta^{-1}$ the objective is dominated by the
observation term and has a benign landscape, and following the solution up to the target keeps the
iterate in that basin. Two solves count as having found the same optimum when the POINTS agree, in
the Hessian metric, and not when their log-densities agree to some number of decimals---in float32
the mode is located only to a few thousandths of a standard deviation, so $\log p$ between two
solves of the same optimum differs by far more than any tight absolute tolerance, and comparing it
reports every restart as a new optimum. On this criterion FitzHugh--Nagumo, Hes1 and HIV each
return a single optimum over five solves. Lorenz returns three, none of them better than the mode,
which is the expected signature of a chaotic objective with shallow secondary basins; the check
passes on the grounds that the mode found is still the best one, and the count is reported so that
the distinction is visible rather than hidden behind the verdict.

\paragraph{Which parameters are identified?} Eigenvalues of the raw Hessian have units, and reading
flatness off them is the same scale error that Section~\ref{sec:solver} is about. We use the
spectrum of $DHD$ and the marginal posterior standard deviation of each parameter as a fraction of
its own value.

\paragraph{Will it be hard to sample?} The Laplace metric whitens the target exactly \emph{at} the
mode, so what matters is how far it departs a standard deviation away. Writing $M = L^\top H(x) L$
for $L$ the mode's whitening, $\operatorname{cond}(M)$ measured over a handful of Laplace draws
predicts the reference chain's behaviour before any chain is run
(Table~\ref{tab:condm}), and costs about eight Hessian evaluations. It is not visible in the
condition number at the mode: on Hes1 the latter is $1.5\times10^5$ while
$\operatorname{cond}(M)$ reaches $10^6$.

\begin{table}[t]
\centering\small
\begin{tabular}{lrrrr}
\toprule
system & $\operatorname{cond}(M)$, Laplace draws & reference draws & reference $\Rhat$ & divergences\\
\midrule
hiv    & \textbf{1.21} & 1.30 & \textbf{1.0001} & \textbf{0\%}\\
lorenz & $1.9\times10^{2}$ & $2.2\times10^{3}$ & 1.0072 & 1.1\%\\
fn     & $3.8\times10^{2}$ & $1.7\times10^{3}$ & 1.0064 & 1.4\%\\
hes1   & $\mathbf{1.1\times10^{6}}$ & $7.1\times10^{6}$ & \textbf{1.76} & \textbf{13\%}\\
\bottomrule
\end{tabular}
\caption{Variation of the curvature over the posterior, and the reference chain that resulted.
Laplace draws reproduce the ordering obtained from reference draws, so the difficulty is knowable
in advance.}
\label{tab:condm}
\end{table}

\paragraph{Is the posterior proper?} A null direction of the quadratic model is not an improper
posterior; the density may still decay at higher order. We walk each parameter axis outward with
the states re-profiled at every step, so the walk follows the ridge rather than cutting across it,
and record the fall in $\log\hat p$. The walk is repeated at several radii and the fall reported at
the furthest one that resolved, since a probe whose inner solve fails also returns $-\infty$ and
would otherwise be indistinguishable from a density that has genuinely decayed. With the
hyperparameters fitted correctly every parameter of every system is proper, and on three of the
four the probe stops resolving before the density measurably decays---so the test now confirms
properness without measuring how strongly, which is weaker than it looks and worth remembering if
it is relied on.

Each check is reported with the threshold it must meet and a pass or fail, alongside an overall
status, so that ``the method declined'' and ``the posterior does not exist'' are not left for the
reader to disentangle. Hes1 is the case where the two reports disagree, and they disagree in the
direction that matters: the \emph{diagnosis} fails, on curvature variation alone---a unique
maximum, no negative curvature, every parameter proper, but $\operatorname{cond}(M) = 1.1\times
10^{6}$---while the \emph{fit} passes its gate, at $\mathrm{ESS} = 21\%$ and $\khat = -0.06$. Both
are right. The profiled integral never sees the funnel, because the states are profiled out and
only seven dimensions are sampled; the reference chain has to traverse it, and cannot. The
practical consequence is uncomfortable and worth stating plainly: on Hes1 we have an estimate that
certifies itself and no way to check it, because the thing we would check it against is the thing
the diagnosis says is infeasible.

One gap in the gate is visible here. On Hes1 $93$ of the $512$ inner profile solves fail outright,
and the gate does not look at that count---it reads the importance weights, which are computed from
the $419$ nodes that succeeded. A high failure rate is not the same as a bad weight distribution
and should probably be gated separately; we report it in the diagnostics but do not currently act
on it.

\begin{table}[t]
\centering\small
\begin{tabular}{llrl}
\toprule
system & parameter & $\mathrm{sd}/|\hat\theta|$ & verdict\\
\midrule
fn     & $a, b, c$ & 0.13, 0.38, 0.02 & identified\\
hiv    & $\lambda, \rho, \delta, N, c$ & 0.04, 0.06, 0.02, 0.02, 0.01 & identified\\
lorenz & $\beta, \rho, \sigma$ & 0.06, 0.03, 0.12 & identified\\
hes1   & $b, c, d, e$ & 0.11 -- 0.26 & identified\\
hes1   & $a, f, g$ & 0.50, 0.56, 0.53 & weak\\
\bottomrule
\end{tabular}
\caption{Identifiability after the hyperparameter fit of Section~\ref{sec:gp}, in five to twelve
seconds per system without sampling. Every system recovers its true parameters to within $1.6$
posterior standard deviations. Before the fix this table read very differently---HIV's $\lambda$
improper, all seven of Hes1's parameters diffuse---and none of that was a property of the models.}
\label{tab:diag}
\end{table}

\section{The profiled posterior}
\label{sec:profiled}

\subsection{Construction}

Do not approximate the states. Integrate them:
\begin{equation}
\label{eq:profile}
  p(\theta) = \int e^{-U(\theta,X)}\,dX
  \;\approx\; e^{-U(\theta, X^\star(\theta))}\,
  \det H_{XX}\big(\theta, X^\star(\theta)\big)^{-1/2},
  \qquad X^\star(\theta) = \arg\min_X U(\theta, X),
\end{equation}
the Laplace approximation applied to the \emph{inner} integral only. The remaining
$p$-dimensional integral is done by self-normalised importance sampling on a scrambled Sobol set,
so the result is deterministic given a seed and comes with an effective sample size and a Pareto
$\khat$ as reference-free diagnostics.

This differs from a Gaussian at the joint mode in three ways that matter. It is exact under
Condition~\ref{cond:A} below, where the joint version still carries a first-order mean error,
because the joint version linearises the $\theta$--$X$ coupling and this one does not. What remains
is the non-Gaussianity of $p(X\mid\theta)$ alone, which the GP prior and the data constrain far
more tightly than they constrain $\theta$. And nothing Gaussian is assumed about $\theta$: skew and
curvature in the parameter marginals survive.

The output is not a Gaussian but a mixture, one Gaussian in the states per parameter node,
\[
  p(x) \;\approx\; \sum_i w_i\, \delta(\theta - \theta_i)\,
  \mathcal N\!\big(X;\, X^\star(\theta_i),\, H_{XX}(\theta_i)^{-1}\big),
\]
from which every moment follows in closed form, including $\operatorname{Cov}(\theta, X)$, which no
Gaussian centred at the mode gets right.

\begin{condition}[Affine in the state]
\label{cond:A}
$f(\cdot,\theta)$ is affine in $X$ for every $\theta$.
\end{condition}

\begin{proposition}
Under Condition~\ref{cond:A} the residual \eqref{eq:residual} is affine in $X$ at fixed $\theta$,
so $U(\theta, X)$ is exactly quadratic in $X$, $p(X\mid\theta)$ is exactly Gaussian, and
\eqref{eq:profile} holds with equality. The $d$-dimensional problem is then exactly a
$p$-dimensional one.
\end{proposition}

Condition~\ref{cond:A} is much weaker than requiring $f$ affine in $(X,\theta)$ jointly---$\theta$
may still multiply states---and covers linear compartment models, linear pharmacokinetics and every
constant-coefficient system. It is testable in two Hessian evaluations by checking whether $H_{XX}$
depends on $X$ \cite{companion}.

\subsection{The proposal, and a structural limit}

Drawing $\theta$ from the joint Laplace marginal $\mathcal N(\hat\theta, (H^{-1})_{\theta\theta})$
is the obvious choice and a poor one: that object accounts for the $\theta$--$X$ coupling only to
first order. At an identical node budget it costs a factor of $28$ in effective sample size on
Lorenz ($67.5\% \to 2.4\%$), $10$ on Hes1 ($21.1\% \to 2.0\%$) and $2.2$ on FitzHugh--Nagumo
($75.7\% \to 35.0\%$). HIV is unaffected ($63.4\% \to 62.6\%$), which is consistent with the rest
of its profile: it is the one system here whose posterior is close to Gaussian in the joint, so
first order is enough. Since $p$ is small, the profiled marginal's own mode and curvature are
directly computable instead---a Newton solve in $p$ dimensions by central differences through
\eqref{eq:profile}, costing $O(p^2)$ profile solves per step.

This is worth doing for the mode. The profiled mode is \emph{not} the joint mode's $\theta$: on
FitzHugh--Nagumo it sits $(-0.24, +1.01, -1.09)$ Laplace standard deviations away, precisely the
directions and magnitudes by which the joint mode is known to be biased. The joint mode maximises
$U(\theta, X)$; the profiled mode maximises $U(\theta, X^\star(\theta)) + \tfrac12\log\det
H_{XX}(\theta)$, and that determinant---the volume of state space consistent with $\theta$---is
what the joint mode ignores.

It is \emph{not} worth doing for the curvature, for a structural reason. The second derivative of
$-U(\theta, X^\star(\theta))$ at the mode is exactly the Schur complement
$(\Sigma_{\theta\theta})^{-1}$, and the determinant term's curvature is negligible: measured on
Hes1 at three stencil sizes spanning two orders of magnitude, the profile curvature equals the
joint Laplace curvature to four decimals. Profiling buys nothing near the mode; everything it gains
is in the shape away from it.

\subsection{A warm start from the implicit function theorem}

$X^\star(\theta)$ satisfies $\nabla_X U(\theta, X^\star) = 0$, so
$H_{XX}\,dX^\star/d\theta + H_{X\theta} = 0$ and the sensitivity is one solve against a
factorisation already in hand. Starting each node from
$X_{\mathrm{MAP}} + (dX^\star/d\theta)(\theta - \theta_{\mathrm{MAP}})$ rather than from
$X_{\mathrm{MAP}}$ buys one halving of the inner iteration count: with the predictor two
iterations already reach the reference floor, without it four are needed
(Table~\ref{tab:predictor}). We had previously reported this as six against two, on the pre-fix
posterior; the effect is real and smaller than that.

\begin{table}[t]
\centering\small
\begin{tabular}{lccc}
\toprule
inner iterations & 2 & 4 & 8\\
\midrule
from $X_{\mathrm{MAP}}$ & 0.0262 & 0.0058 & 0.0121\\
\textbf{from the predictor} & \textbf{0.0120} & 0.0127 & 0.0128\\
\bottomrule
\end{tabular}
\caption{Largest parameter error in reference standard deviations on FitzHugh--Nagumo, against a
chain-to-chain floor of $0.0100$. Everything at or below about $0.013$ is at the floor and the
differences among those entries are Monte Carlo noise; the one entry that is clearly off it is
$X_{\mathrm{MAP}}$ at two iterations. The predictor's benefit is that two iterations suffice, which
is roughly a quarter of the wall clock.}
\label{tab:predictor}
\end{table}

\subsection{Choosing the finite-difference step}
\label{sec:stencil}

The Newton step of the previous subsection differentiates \eqref{eq:profile} by central
differences, and the step size $h$ is not a free parameter to be set once. A central second
difference carries truncation error growing like $h^2$ times the fourth derivative and round-off
error falling like $\sigma/h^2$, so $h$ is bounded on both sides, and \emph{both bounds are
properties of the problem}: $\sigma$ comes from the working precision, the fourth derivative from
how non-Gaussian the profiled marginal is. Table~\ref{tab:stencil} shows the two bounds measured
across the suite, in units of the joint Laplace standard deviation.

\begin{table}[t]
\centering\small
\begin{tabular}{llrrrrrrr}
\toprule
& $h/\mathrm{sd}$ & 0.05 & 0.2 & 0.8 & 1.6 & 3.2 & 6.4 & 12.8\\
\midrule
fn, float32   & $\theta_2$ & \textbf{--0.098} & 1.028 & 1.063 & 1.068 & 1.073 & 1.086 & 1.126\\
fn, float64   & $\theta_1$ & 1.112 & 1.112 & 1.116 & 1.128 & 1.176 & 1.359 & \textbf{1.967}\\
hiv, float64  & $\theta_0$ & 21.330 & 21.330 & 21.330 & 21.330 & 21.330 & 21.330 & 21.330\\
lorenz, f64   & $\theta_2$ & 1.056 & 1.058 & 1.076 & 1.138 & 1.441 & \textbf{4.198} & 1.191\\
hes1, float64 & $\theta_5$ & 216.9 & 215.7 & \textbf{192.1} & --- & --- & --- & ---\\
\bottomrule
\end{tabular}
\caption{Diagonal curvature of the profiled log-density, scaled so a Gaussian profile would read
$1.0$ at every $h$, after the hyperparameter fit of Section~\ref{sec:gp}. Four regimes appear.
Round-off at small $h$ in single precision drives FitzHugh--Nagumo's $\theta_2$ \emph{negative}.
Truncation bends the estimate at large $h$, gently on FitzHugh--Nagumo in double and sharply on
Hes1, whose $\theta_5$ has already moved $12\%$ by $h = 0.8$ and whose probes stop resolving
beyond $1.6$. Lorenz shows both, and the non-monotone excursion at $6.4$ is a reminder that a
two-point comparison can land on either side of a spike. HIV is flat to five figures across three
orders of magnitude, which is why its plateau runs to the top of the ladder.}
\label{tab:stencil}
\end{table}

Since a plateau is precisely a range over which the estimate does not move, we locate it: evaluate
the diagonal curvature along a ladder, take the longest \emph{contiguous run of rungs that all
agree with one another} to within a tolerance, and use the geometric middle of that run. Three
details are each necessary, and we record them because the obvious version of the rule fails on
this suite in three different ways.

\emph{Agreement must be across the run, not between neighbours.} On FitzHugh--Nagumo in double
precision every consecutive change stays under $5\%$ from $h = 0.05$ out to $3.2$ while the
curvature drifts from $1.112$ to $1.176$; small steps accumulate, and a neighbour test walks
steadily off the plateau it is meant to find.

\emph{The middle of the run, not the top.} The ends of the run are exactly where the estimate
begins to fail, and both bounds are multiplicative, so the geometric middle maximises the margin
on either side. It matters too that the run is located from diagonal probes while the Newton also
needs off-diagonal ones at $\sqrt2$ times the displacement, so a step only just inside the run for
the first can be outside it for the second. This one we must qualify: the measurement behind it was
Lorenz \emph{before} the hyperparameter fit, where taking the top of the run cost three quarters of
the effective sample size, and it no longer reproduces---Lorenz's run now ends at $0.8$ and taking
the top there costs nothing ($67.8\%$ against $67.5\%$). We keep the choice because it is free and
cannot hurt, not because we can still exhibit a case that needs it.

\emph{A failed probe is not agreement.} Judging coherence on whichever parameters happen to be
finite silently drops those that failed, and the parameters do not fail together. Hes1 is the live
case: at $h = 1.6$ its $\theta_5$ probe has already stopped resolving while $\theta_2$, $\theta_3$
and $\theta_4$ are still moving smoothly, and at $3.2$ three more have gone while those three
continue. Judged on the survivors the plateau would run to $6.4$, an order of magnitude past the
point where the parameter with the sharpest profile stopped being measurable at all. A rung counts
only if every parameter resolved on it.

The cost is $2p$ evaluations per rung, negligible beside the node budget. The selected steps differ
by more than an order of magnitude across the suite (Table~\ref{tab:chosen}), which is the case for
selecting rather than setting them. The margin available to a constant is thinner than that number
suggests: FitzHugh--Nagumo in single precision needs $h \gtrsim 0.4$ before round-off stops
dominating, and Hes1 needs $h \lesssim 0.4$ before truncation takes over, so the two windows meet
at a point. A fixed step would have to be that point, and would sit on the edge of both.

\subsection{A gate}

Importance sampling at a few percent effective sample size is not an estimate. We require
$\mathrm{ESS} \ge 10\%$ and $\khat < 0.7$, both computable without a reference, and otherwise
report the Laplace approximation.

We should be candid that the gate no longer fires on this suite in double precision: after the
hyperparameter fit all four systems pass it, Hes1 least comfortably at $21\%$. Its one live decline is Hes1 in float32, at
$\mathrm{ESS} = 2.1\%$. The case that used to justify it---a
posterior where nothing was identified, the mode already sat at the reference mean, and profiling
made matters worse---was an artifact of the broken hyperparameters and no longer exists. We keep
the gate because the failure it guards against is real and cheap to check for, not because we can
still exhibit it; a reader should weigh it as a safeguard of unproven necessity rather than as a
demonstrated one.

\begin{algorithm}[t]
\DontPrintSemicolon
\SetAlgoLined
\KwIn{MAGI model; nodes $N=512$; inner iterations $3$; gate $(\mathrm{ESS}\ge10\%,\ \khat<0.7)$}
\KwOut{profiled mixture, or Laplace fallback, with diagnostics}
$\xmap \leftarrow$ Jacobi-scaled Gauss--Newton on $\|R\|^2$ \tcp*[r]{$\approx 1$\,s}
$H \leftarrow J^\top J + \sum_a R_a\nabla^2R_a$; Jacobi-stabilised $\Sigma \leftarrow H^{-1}$\;
$dX^\star/d\theta \leftarrow -H_{XX}^{-1}H_{X\theta}$ \tcp*[r]{implicit function theorem}
$h \leftarrow$ geometric middle of the curvature plateau \tcp*[r]{Sec.~\ref{sec:stencil}}
$\hat\theta, H_{\mathrm{prof}} \leftarrow$ Newton on $\log\hat p(\theta)$, central differences at $h$\;
$\theta_i \leftarrow \hat\theta + L z_i$, $z_i$ scrambled Sobol, $LL^\top = H_{\mathrm{prof}}^{-1}$\;
\ForEach{$\theta_i$ \rm{(batched)}}{
  $X \leftarrow X_{\mathrm{MAP}} + (dX^\star/d\theta)(\theta_i - \theta_{\mathrm{MAP}})$\;
  three Gauss--Newton steps on $X$ at fixed $\theta_i$\;
  $\log\hat p_i \leftarrow -U(\theta_i, X) - \tfrac12\log\det H_{XX}$\;
}
$w \propto \exp(\log\hat p_i - \log q_i)$; compute ESS and $\khat$\;
\Return{mixture if the gate passes, else $\mathcal N(\xmap, \Sigma)$}
\caption{Profiled MAGI posterior, $\approx 5$\,s at $d=325$}
\label{alg:profiled}
\end{algorithm}

\section{Benchmarks}
\label{sec:bench}

\subsection{Systems, and a note on constructing them}

Four systems: \textbf{FitzHugh--Nagumo} ($d=325$, $p=3$), \textbf{Hes1} ($d=106$, $p=7$, a
rational nonlinearity, $\exp(X)$ states and a component never observed), \textbf{HIV} ($d=608$,
$p=5$, time-dependent forcing, states spanning $30$ to $10^5$) and \textbf{Lorenz} ($d=306$,
$p=3$, chaotic).

The suite originally integrated with forward Euler, retaining every substep. Both choices bound
what could be measured. Storage is $(t_{\max}/h)\,D$, which at $h=10^{-6}$ is $5.8$\,GB for Hes1's
horizon of $240$; and on HIV the integration error at the observation times was $2.6$ observation
standard deviations, larger than the noise it was competing with, so the posterior under study was
not the intended one. Replacing Euler with RK4 and retaining only the output grid changes both:

\begin{center}\small
\begin{tabular}{lrrrr}
\toprule
& \multicolumn{2}{c}{error / $\sigma$} & \multicolumn{2}{c}{$f$ evaluations}\\
\cmidrule(lr){2-3}\cmidrule(lr){4-5}
system & Euler $10^{-6}$ & \textbf{RK4 $10^{-3}$} & Euler $10^{-6}$ & \textbf{RK4 $10^{-3}$}\\
\midrule
fn     & 6.8e-05 & \textbf{6.8e-11} & 20{,}000{,}000 & \textbf{80{,}000}\\
hes1   & 2.2e-06 & \textbf{7.5e-12} & 240{,}000{,}000 & \textbf{960{,}000}\\
hiv    & 5.0e-03 & \textbf{5.2e-09} & 20{,}000{,}200 & \textbf{80{,}800}\\
lorenz & 1.7e-04 & \textbf{1.1e-08} & 2{,}500{,}100 & \textbf{10{,}400}\\
\bottomrule
\end{tabular}
\end{center}

Six orders of accuracy for two to three fewer of work, with storage independent of the step: $1$ to
$5$\,kB rather than gigabytes. Hes1's integration falls from $10.0$\,s to $0.14$\,s. We record this
because a benchmark whose synthetic ``truth'' carries more error than its own noise cannot
discriminate between posterior approximations, and the defect is invisible unless looked for.

\subsection{References}

References are $32$ chains of $3{,}000$ draws of NUTS \cite{hoffman2014} via BlackJAX,
preconditioned by the exact Hessian in Jacobi-scaled coordinates, which affects only efficiency and
never the stationary distribution. Rebuilt after the hyperparameter fit, HIV gives
$\Rhat = 1.0001$ with \emph{zero} divergences---it had previously failed to converge at all---
FitzHugh--Nagumo $1.0064$ with $1.4\%$ and Lorenz $1.0072$ with $1.1\%$. Hes1 gives
$\Rhat = 1.76$ with $13\%$ divergences and $2.79$ million leapfrog steps per chain, and is not
usable; it was marginally usable \emph{before} the fix, for reasons Section~\ref{sec:funnel}
sets out.

\subsection{Results}
\label{sec:results}

\begin{table}[t]
\centering\small
\begin{tabular}{llrrr}
\toprule
system & estimate & max $|\theta$ error$|$ & max $|$sd error$|$ & s\\
\midrule
fn & mode                    & 1.0286 & 2.18\% & ---\\
   & \textbf{profiled}       & \textbf{0.0126} & \textbf{0.45\%} & 6.7\\
   & reference half-vs-half  & 0.0100 & 0.86\% & \\
\midrule
hiv & mode                   & 0.1485 & 0.62\% & ---\\
   & \textbf{profiled}       & \textbf{0.0093} & \textbf{0.34\%} & 23.2\\
   & reference half-vs-half  & 0.0081 & 0.77\% & \\
\midrule
lorenz & mode                & 1.8026 & 39.68\% & ---\\
   & \textbf{profiled}       & \textbf{0.0119} & \textbf{1.27\%} & 5.2\\
   & reference half-vs-half  & 0.0405 & 1.19\% & \\
\midrule
hes1$^\dagger$ & mode        & 2.8639 & 189.6\% & ---\\
   & profiled                & 0.2163 & 50.5\% & 3.2\\
   & reference half-vs-half  & 0.0841 & 10.9\% & \\
\bottomrule
\end{tabular}
\caption{Parameter errors in reference standard deviations, after the hyperparameter fit of
Section~\ref{sec:gp}. On FitzHugh--Nagumo, HIV and Lorenz the profiled posterior is at or
\emph{below} the level at which two independent halves of the reference agree with each other.
Lorenz in particular moves from $0.34$ before the fix---which we had attributed to chaotic dynamics
defeating the inner Laplace approximation---to $0.0119$, so that explanation was wrong.
$^\dagger$Hes1's reference has $\Rhat = 1.76$ and $13\%$ divergences and is not usable; its row is
shown for completeness only.}
\label{tab:results}
\end{table}

Table~\ref{tab:results} is the main result. On all three systems with a usable reference the
profiled posterior sits at or below the level at which two independent halves of that reference
agree with each other, which is the most that can be claimed from a finite chain: $0.0126$ against
a floor of $0.0100$ on FitzHugh--Nagumo, $0.0093$ against $0.0081$ on HIV, and $0.0119$ against
$0.0405$ on Lorenz, the last comfortably \emph{below} it. Against the mode the improvement is $82$,
$16$ and $151$ times on location, and the spread errors fall to a fraction of a percent on the
first two. The costs in the last column---between three and twenty-three seconds---are against
reference chains of roughly two hundred seconds and up.

Lorenz deserves a note, because we previously drew a conclusion from it that was wrong. Before the
hyperparameter fit its profiled error was $0.34$ where the mode's was $4.14$, and we attributed the
gap to chaotic dynamics making $p(X\mid\theta)$ strongly non-Gaussian and so undermining the inner
Laplace approximation \eqref{eq:profile}. That was a plausible mechanism and it was not the one
operating: with $\phi$ fitted properly the error is $0.0119$, below the floor, and the chaos costs
nothing measurable. Whether \eqref{eq:profile} degrades on genuinely chaotic problems remains open;
this suite no longer speaks to it.

Hes1's row is shown for completeness and should not be read as a result. Its gate passes
($\mathrm{ESS} = 21\%$), but its reference does not converge, so $0.2163$ is a comparison against a
chain that has not found the posterior rather than a measurement of error. Section~\ref{sec:funnel}
takes up why.

\subsection{Precision, and what the stencil was hiding}

The profiled estimate appeared at first to require double precision: in single precision its
effective sample size on FitzHugh--Nagumo fell from $75.7\%$ to $2.4\%$ and the gate replaced it
with the very approximation it was meant to improve on. That diagnosis was wrong, and the way it
was wrong is worth recording.

The importance weights are barely affected by working precision. Measuring the float32 error in
$\log\hat p$ across profile nodes---and only its \emph{spread} matters, since a constant offset
cancels in self-normalisation---gives $0.007$ nats on CPU and $0.052$ on GPU, which degrade the
effective sample size by a factor of $0.997$. Guarding the one unguarded matrix--matrix product in
the exact-Hessian assembly, which runs in TF32 by default on tensor-core hardware, improves the
log-determinant's error spread by a factor of ten on GPU and changes the effective sample size not
at all.

What was actually failing was the finite-difference stencil. A central second difference divides by
$h^2$, so at $h = 0.02\,\mathrm{sd}$ the $0.007$ nats of round-off become an amplification of
roughly $2500$ and the curvature estimate is noise---on FitzHugh--Nagumo it comes out
\emph{negative}. With the step chosen as in Section~\ref{sec:stencil} rather than fixed, single
precision matches double on FitzHugh--Nagumo and Lorenz and beats it on HIV, in every case by
selecting a step an octave or two larger to stay clear of its own round-off floor
(Table~\ref{tab:chosen}).

The exception is Hes1, and it is a real one: single precision takes its effective sample size from
$21\%$ to $2.1\%$ and the gate declines. It is not the stencil---both precisions pick a step inside
the same plateau---and we are not able to attribute it further. Hes1 is the system with the funnel
(Section~\ref{sec:funnel}), so the curvature the inner solves and the log-determinant see varies by
six orders of magnitude across the nodes, and float32 has less margin to absorb that; but HIV's
Hessian is worse conditioned at the mode and single precision helps it, so conditioning alone is
not the story and we have not isolated what is. The honest summary is narrower than we first wrote
it: use single precision, and expect the gate to catch the case where it is not enough.

\begin{table}[t]
\centering\small
\begin{tabular}{lrrrrrr}
\toprule
system & \multicolumn{2}{c}{chosen $h/\mathrm{sd}$} & \multicolumn{2}{c}{plateau} & \multicolumn{2}{c}{ESS}\\
& f64 & f32 & f64 & f32 & f64 & f32\\
\midrule
fn     & 0.1  & 0.8  & 0.01--1.6  & 0.2--1.6   & 75.7\% & 76.0\%\\
hiv    & 0.2  & 1.6  & 0.01--6.4$^\ast$ & 0.4--6.4$^\ast$ & 63.4\% & 71.8\%\\
lorenz & 0.1  & 0.4  & 0.01--0.8  & 0.2--0.8   & 67.5\% & 67.9\%\\
hes1   & 0.05 & 0.1  & 0.01--0.4  & 0.05--0.4  & 21.1\% & \textbf{2.1\%}\\
\bottomrule
\end{tabular}
\caption{Steps selected automatically, and the plateaus they were selected from, after the
hyperparameter fit of Section~\ref{sec:gp}. The lower edge moves with precision---FitzHugh--Nagumo
$0.01$ in double against $0.4$ in single---because round-off sets it; the upper edge does not
(Lorenz at $0.8$, Hes1 at $0.4$) because truncation and probe failure set it. The rule reads both
without being told which regime it is in. $^\ast$The plateau runs to the top of the ladder, so its
upper edge is the ladder's and not the problem's; the step taken is the geometric middle either
way.}
\label{tab:chosen}
\end{table}

We note in passing a fact visible only in these diagnostics: Lorenz's profiled curvature for
$\sigma$ is half the Laplace value, and for $\rho$ four fifths, across the whole plateau. The
Laplace covariance is materially wrong there in a way the curvature at the mode cannot reveal,
which is the effect profiling exists to capture.

\section{Negative results}
\label{sec:negative}

\emph{The third-order correction of \cite{companion} is superseded.} It helped on FitzHugh--Nagumo
and harmed Lorenz ($4.14 \to 4.86$) and Hes1. Those figures predate the hyperparameter fit and so
describe different posteriors; we did not re-measure them, because the profiled estimate is better
than either under both.

\emph{Correcting the inner Laplace approximation by sampling makes matters worse.} Writing the
exact inner integral as $-U^\star - \tfrac12\log\det H_{XX} + \log \E_\xi[e^{-\Delta}]$ with
$\xi \sim \mathcal N(0, H_{XX}^{-1})$, the omitted term can be estimated with antithetic pairs,
which cancel $\Delta$'s cubic part exactly. It fails: $\Delta$ is a sum of quartic terms over $322$
to $603$ coordinates, a handful of pairs cannot resolve its mean, and injecting that noise into
$\log\hat p(\theta)$ corrupts the weights multiplicatively. FitzHugh--Nagumo goes from $0.0086$ to
$0.3665$. That it reaches the reference floor with the correction \emph{off} is the useful part.

\emph{Restricting the integration to an identifiable subspace did not rescue HIV}, when HIV still
appeared to need rescuing. Profiling the apparently improper directions out rather than integrating
them is well defined---the determinant becomes a pseudo-determinant over the complement---but the
effective sample size stayed at $1.6\%$ and half the inner solves failed. We record it because the
construction is sound and may be useful on a genuinely partially identified problem; the target it
was built for turned out to be an artifact.

\emph{Profiling buys nothing near the mode}, for the structural reason in
Section~\ref{sec:profiled}.

\emph{Two convergence tests that report success on failure.} Clamping the Newton decrement with
$\max(g^\top\delta, 0)$, to guard against a negative value, converts a round-off failure into a
reported convergence: on HIV in float32 the dot product comes out negative, the clamp returns
zero, and the solver stops after one iteration at a point $246$ nats below the mode. A
non-positive or non-finite decrement must be treated as \emph{not converged}. Separately, counting
distinct optima by rounding $\log p$ reports every restart as a new optimum in single precision.
Both produced confident, wrong answers rather than visible failures.

\emph{A fixed finite-difference step is not safe}, and the three natural ways of choosing one
adaptively each fail on this suite: a neighbour-agreement test walks off the plateau, the top of
the plateau breaks the off-diagonal probes, and treating a failed probe as agreement extends the
plateau past its real end. Section~\ref{sec:stencil} gives the measurements. We record these
because the first two produced plausible-looking answers on three of the four systems and failed
only on the fourth.

\section{Limitations}

The evidence is four systems from one benchmark suite, with usable references on three of them.
The profiled posterior reaches the reference floor on all three, which is as much as a finite
reference can establish and no more: it bounds the error by the chain's own Monte Carlo noise
rather than measuring it. A sharper test would need either a longer reference or a problem with a
known answer.

We also no longer have a case in this suite where the inner Laplace approximation
\eqref{eq:profile} visibly degrades. We thought Lorenz was one and it was not
(Section~\ref{sec:results}); Condition~\ref{cond:A} remains the right diagnostic for it, and
remains unrun against a case that fails.

The gate's thresholds ($\mathrm{ESS} \ge 10\%$, $\khat < 0.7$) are conventional rather than
calibrated here, and were correct on all four systems, which is weak evidence.

Hes1's reference is unusable, so on one system of four our comparisons have no target;
Section~\ref{sec:funnel} sets out why and what might be done about it.

Finally, we argued that the choice of prior on unidentified parameters determines the estimates of
the identified ones. The argument stands, but our evidence for it does not: the example was HIV,
which is not in fact unidentified. Whether the effect matters in practice is therefore open, and
would need a problem that is genuinely partially identified once its hyperparameters are fitted
properly.

\section{A funnel, and how one might sample it}
\label{sec:funnel}

Hes1 is the one system in the suite without a usable reference, and the reason is worth setting out
because it is diagnosable and because the remedies are not the obvious ones.

\subsection{What it is}

Its local curvature varies by a factor of $35.7$ across its own posterior, against $1.01$ for HIV.
Regressing $\log$ of the largest Hessian eigenvalue on the coordinates over reference draws
identifies the drivers: the peak of the third state component, correlation $+0.933$, and the
parameter $a$, correlation $-0.902$ with a slope of $-97.6$ per unit. Hes1's states are on a log
scale, its field reading $P, M, H = \exp(X)$, and $a$ enters as $-aH$ and $-aP$. So a parameter
multiplies an exponentiated state, and the state in question, $H$, is \emph{never observed}. That
is a funnel in the textbook sense, and the coordinate that generates it carries no data.

\subsection{Why it appeared only after the hyperparameter fix}

Before the fix Hes1's parameters were unidentified and pinned near zero; with $a \approx 10^{-9}$
the $aH$ coupling was inert, the posterior was flat, and NUTS wandered it comfortably to
$\Rhat = 1.005$ while exploring a region that meant nothing. Afterwards $a$ is determined at
$0.032$, the coupling is live, and the same sampler takes $2.79$ million leapfrog steps per chain
and reaches $\Rhat = 1.76$. This is not a regression: the problem went from flat, easy and
meaningless to concentrated, hard and correct.

\subsection{No fixed metric can do it}

The obstruction is not the conditioning at any one point but the mutual disagreement between
Hessians at different points, since a mass matrix must whiten all of them at once:

\begin{center}\small
\begin{tabular}{lrrr}
\toprule
system & pairwise $\operatorname{cond}(H_i^{-1}H_j)$ & mode metric & averaged metric\\
\midrule
hiv    & 1.18 & 1.18 & 1.14\\
fn     & $2.1\times10^{3}$ & $3.8\times10^{2}$ & 59\\
lorenz & $9.5\times10^{10}$ & $1.1\times10^{2}$ & 46\\
hes1   & $\mathbf{1.7\times10^{13}}$ & $3.7\times10^{5}$ & $\mathbf{5.2\times10^{11}}$\\
\bottomrule
\end{tabular}
\end{center}

On Hes1 the averaged metric is six orders \emph{worse} than the mode's, so the cheap alternatives
are ruled out without trying them. Lorenz is the instructive near-miss: its pairwise figure is also
enormous, yet a good fixed metric still reduces it to $46$.

\subsection{What we would try next}

\begin{enumerate}[leftmargin=1.5em,itemsep=2pt]
\item \textbf{Observe $H$.} The funnel is generated by a component with no data. If the
  experiment admits any observation of it the geometry largely dissolves. This is a change to the
  problem rather than to the method, and it is by some distance the cheapest item here; it is worth
  establishing before building machinery.
\item \textbf{Reparameterise to the profiled coordinates.} Sampling
  $(\theta, y)$ with $X = X^\star(\theta) + L(\theta)y$ removes the $\theta$--$X$ coupling by
  construction, and is exactly the decomposition Section~\ref{sec:profiled} already computes. The
  obstacle is that HMC needs gradients through an argmin and a log-determinant; the implicit
  function theorem supplies the first, and the second requires third derivatives of $U$, which the
  structure of \eqref{eq:residual} makes available---$\nabla^2 R_a = \sqrt b (L_k^\top)_a\nabla^2f$,
  so a third derivative is one more level on the user's vector field alone.
\item \textbf{Riemannian HMC with a SoftAbs metric}, the standard answer to a position-dependent
  geometry. It needs the same third derivatives and is slow per step, but it does not require the
  profiled construction to be correct.
\item \textbf{MCMC on the profiled parameter density.} With $p = 7$ a derivative-free sampler is
  entirely practical and would settle whether the importance sampling of
  Section~\ref{sec:profiled} is adequate. It is worth being clear about what this does not do: it
  shares the inner Laplace approximation with the method under test, so it validates the
  parameter-space integration and nothing else. It is not an independent reference.
\end{enumerate}

We note that the profiled posterior is structurally unaffected by this particular funnel, since it
integrates the states out at each $\theta$ instead of traversing the joint geometry; on Hes1 it
reports an effective sample size of $21\%$ and $\khat = -0.06$ where NUTS fails outright. That is
suggestive rather than established---without a reference we cannot score the answer---but if it
holds it is an argument for the method that has nothing to do with speed.

\section{Conclusion}

MAGI's posterior is mostly state and a little parameter, and the two deserve different treatment.
Integrating the states out exactly at each parameter value, and doing the small remaining integral
directly, reaches the accuracy of a long reference chain in seconds and reports whether it has.
Getting there depended on four unglamorous corrections---scaling a linear solve, integrating a
benchmark accurately, restoring a prior term the derivation always contained, and fitting the
Gaussian process hyperparameters on a scale-free footing---each of which was invisible in the
outputs until it was looked for. The last was the largest by some distance, and the most
instructive: it did not make the method fail, it made it succeed at fitting the wrong posterior,
and every certificate we had built agreed that the fit was good. What caught it was plotting the
trajectories.

A pattern runs through all of them. Six separate quantities in this pipeline were being compared
against absolute thresholds while carrying units: the conditioning of the normal equations, the
eigenvalues used to detect flat directions, the finite-difference step, the convergence test, the
agreement test between restarts, and the starting point of the hyperparameter optimiser. Each
worked on the problem it was developed on and failed on another, and in every case the repair was
the same---divide by the scale that the quantity is measured in, and the threshold becomes a
statement about the posterior rather than about the units. It is a dull principle, and it accounted
for more of the difficulty here than any of the approximation theory. The last of the six is the
one to remember, because it did not announce itself as a numerical problem at all: it produced a
posterior that was internally consistent, passed every certificate we had, and answered a different
question than the one asked.

The seventh constant, the ridge before a Cholesky, is worth keeping beside them as a warning
against applying the principle mechanically. It looks identical---a tolerance calibrated in one
precision, used in another---and scaling it to the dtype fixed the symptom and silently replaced
the objective. The distinction is whether the number is being \emph{compared} against something or
\emph{added} to it. Only the first kind can be made scale-free by dividing.

\begin{thebibliography}{9}
\small
\bibitem{yang2021} S.~Yang, S.~W.~K.~Wong, S.~C.~Kou.
  Inference of dynamic systems from noisy and sparse data via manifold-constrained Gaussian
  processes. \emph{PNAS} 118(15), 2021.
\bibitem{companion} A fast, self-certifying Gaussian posterior for MAGI ODE parameter inference.
  Companion technical report, 2026.
\bibitem{tierney1986} L.~Tierney, J.~B.~Kadane.
  Accurate approximations for posterior moments and marginal densities.
  \emph{JASA} 81(393):82--86, 1986.
\bibitem{hoffman2014} M.~D.~Hoffman, A.~Gelman.
  The No-U-Turn sampler. \emph{JMLR} 15:1593--1623, 2014.
\bibitem{vehtari2024} A.~Vehtari, D.~Simpson, A.~Gelman, Y.~Yao, J.~Gabry.
  Pareto smoothed importance sampling. \emph{JMLR} 25:1--58, 2024.
\bibitem{marquardt1963} D.~W.~Marquardt.
  An algorithm for least-squares estimation of nonlinear parameters.
  \emph{J.~SIAM} 11(2):431--441, 1963.
\end{thebibliography}

\end{document}
